%=========================================================================
%  Séance 0 : Présentation du cours
%  Compilation :  pdflatex 00-presentation.tex   (deux fois)
%=========================================================================
\documentclass[10pt,aspectratio=169]{beamer}
%=========================================================================
%  Préambule commun — Analyse des données, L3 Économie
%  Appelé par chaque séance via %=========================================================================
%  Préambule commun — Analyse des données, L3 Économie
%  Appelé par chaque séance via %=========================================================================
%  Préambule commun — Analyse des données, L3 Économie
%  Appelé par chaque séance via \input{preambule-cy}
%=========================================================================
\usetheme[progressbar=frametitle,numbering=fraction,block=fill]{metropolis}

\usepackage[T1]{fontenc}
\usepackage[utf8]{inputenc}
\usepackage[french]{babel}
%\usepackage{lmodern}  % absent de cet environnement de compilation ; a decommenter
\usepackage{booktabs}
\usepackage{listings}
\usepackage{xcolor}

%-------------------------------------------------------------------------
%  PALETTE — remplacer par les valeurs de votre charte CY habituelle
%-------------------------------------------------------------------------
\definecolor{CYprincipal}{HTML}{1B2A4A}   % couleur dominante (titres, barres)
\definecolor{CYaccent}{HTML}{C8102E}      % couleur d'accentuation
\definecolor{CYgris}{HTML}{55637C}        % texte secondaire
\definecolor{CYclair}{HTML}{EEF1F6}       % fonds de blocs

\setbeamercolor{palette primary}{bg=CYprincipal,fg=white}
\setbeamercolor{progress bar}{fg=CYaccent}
\setbeamercolor{frametitle}{bg=CYprincipal,fg=white}
\setbeamercolor{alerted text}{fg=CYaccent}
\setbeamercolor{block title}{bg=CYclair,fg=CYprincipal}
\setbeamercolor{block body}{bg=CYclair!60,fg=black}

%-------------------------------------------------------------------------
%  Mise en forme du code Python
%-------------------------------------------------------------------------
\lstdefinestyle{py}{
  language=Python,
  basicstyle=\ttfamily\small,
  keywordstyle=\color{CYaccent},
  stringstyle=\color{CYprincipal},
  commentstyle=\color{CYgris}\itshape,
  showstringspaces=false,
  breaklines=true,
  frame=none,
  literate={é}{{\'e}}1 {è}{{\`e}}1 {à}{{\`a}}1 {ç}{{\c c}}1 {ê}{{\^e}}1 {ô}{{\^o}}1 {û}{{\^u}}1 {î}{{\^i}}1
}
\lstset{style=py}

\author{Stefania Marcassa}
\institute{CY Cergy Paris Université --- L3 Économie}
\date{}

%=========================================================================
\usetheme[progressbar=frametitle,numbering=fraction,block=fill]{metropolis}

\usepackage[T1]{fontenc}
\usepackage[utf8]{inputenc}
\usepackage[french]{babel}
%\usepackage{lmodern}  % absent de cet environnement de compilation ; a decommenter
\usepackage{booktabs}
\usepackage{listings}
\usepackage{xcolor}

%-------------------------------------------------------------------------
%  PALETTE — remplacer par les valeurs de votre charte CY habituelle
%-------------------------------------------------------------------------
\definecolor{CYprincipal}{HTML}{1B2A4A}   % couleur dominante (titres, barres)
\definecolor{CYaccent}{HTML}{C8102E}      % couleur d'accentuation
\definecolor{CYgris}{HTML}{55637C}        % texte secondaire
\definecolor{CYclair}{HTML}{EEF1F6}       % fonds de blocs

\setbeamercolor{palette primary}{bg=CYprincipal,fg=white}
\setbeamercolor{progress bar}{fg=CYaccent}
\setbeamercolor{frametitle}{bg=CYprincipal,fg=white}
\setbeamercolor{alerted text}{fg=CYaccent}
\setbeamercolor{block title}{bg=CYclair,fg=CYprincipal}
\setbeamercolor{block body}{bg=CYclair!60,fg=black}

%-------------------------------------------------------------------------
%  Mise en forme du code Python
%-------------------------------------------------------------------------
\lstdefinestyle{py}{
  language=Python,
  basicstyle=\ttfamily\small,
  keywordstyle=\color{CYaccent},
  stringstyle=\color{CYprincipal},
  commentstyle=\color{CYgris}\itshape,
  showstringspaces=false,
  breaklines=true,
  frame=none,
  literate={é}{{\'e}}1 {è}{{\`e}}1 {à}{{\`a}}1 {ç}{{\c c}}1 {ê}{{\^e}}1 {ô}{{\^o}}1 {û}{{\^u}}1 {î}{{\^i}}1
}
\lstset{style=py}

\author{Stefania Marcassa}
\institute{CY Cergy Paris Université --- L3 Économie}
\date{}

%=========================================================================
\usetheme[progressbar=frametitle,numbering=fraction,block=fill]{metropolis}

\usepackage[T1]{fontenc}
\usepackage[utf8]{inputenc}
\usepackage[french]{babel}
%\usepackage{lmodern}  % absent de cet environnement de compilation ; a decommenter
\usepackage{booktabs}
\usepackage{listings}
\usepackage{xcolor}

%-------------------------------------------------------------------------
%  PALETTE — remplacer par les valeurs de votre charte CY habituelle
%-------------------------------------------------------------------------
\definecolor{CYprincipal}{HTML}{1B2A4A}   % couleur dominante (titres, barres)
\definecolor{CYaccent}{HTML}{C8102E}      % couleur d'accentuation
\definecolor{CYgris}{HTML}{55637C}        % texte secondaire
\definecolor{CYclair}{HTML}{EEF1F6}       % fonds de blocs

\setbeamercolor{palette primary}{bg=CYprincipal,fg=white}
\setbeamercolor{progress bar}{fg=CYaccent}
\setbeamercolor{frametitle}{bg=CYprincipal,fg=white}
\setbeamercolor{alerted text}{fg=CYaccent}
\setbeamercolor{block title}{bg=CYclair,fg=CYprincipal}
\setbeamercolor{block body}{bg=CYclair!60,fg=black}

%-------------------------------------------------------------------------
%  Mise en forme du code Python
%-------------------------------------------------------------------------
\lstdefinestyle{py}{
  language=Python,
  basicstyle=\ttfamily\small,
  keywordstyle=\color{CYaccent},
  stringstyle=\color{CYprincipal},
  commentstyle=\color{CYgris}\itshape,
  showstringspaces=false,
  breaklines=true,
  frame=none,
  literate={é}{{\'e}}1 {è}{{\`e}}1 {à}{{\`a}}1 {ç}{{\c c}}1 {ê}{{\^e}}1 {ô}{{\^o}}1 {û}{{\^u}}1 {î}{{\^i}}1
}
\lstset{style=py}

\author{Stefania Marcassa}
\institute{CY Cergy Paris Université --- L3 Économie}
\date{}


\title{Présentation du cours}
\subtitle{Analyse des données --- Séance 0}

\begin{document}

\maketitle

%=========================================================================
\section{Ce que nous allons faire}
%=========================================================================

\begin{frame}{Ce que vous saurez faire en décembre}
\begin{enumerate}
  \item \textbf{Trouver et charger} des données économiques réelles, et lire la documentation d'une source.
  \item \textbf{Nettoyer et restructurer} un jeu de données brut.
  \item \textbf{Produire des statistiques descriptives défendables}.
  \item \textbf{Construire des graphiques honnêtes}.
  \item \textbf{Estimer et lire une régression} en pratique.
  \item \textbf{Rendre un travail reproductible}.
\end{enumerate}
\vfill
\begin{alertblock}{Le critère}
Qu'une autre personne puisse ré-exécuter votre code et retrouver vos résultats.
\end{alertblock}
\end{frame}

\begin{frame}{Ce que ce cours n'est pas}
Ce n'est pas un cours d'économétrie théorique. Les propriétés des estimateurs sont supposées connues.

\vspace{1em}
C'est un cours de \alert{mise en œuvre} : la difficulté est dans les données, pas dans les formules.

\vspace{1.5em}
\begin{block}{Une mise en garde, dès aujourd'hui}
Une régression ne mesure pas un effet causal. Elle mesure une association conditionnelle aux variables incluses.

\vspace{0.4em}
Tout le métier tient dans l'écart entre les deux. Nous y reviendrons à partir de la séance~7.
\end{block}
\end{frame}

\begin{frame}{Les dix séances}
\small
\begin{tabular}{clp{6.2cm}}
\toprule
& \textbf{Séance} & \textbf{Application en TD} \\
\midrule
1 & Environnement et anatomie des données & Indice des prix \\
2 & Chargement, inspection, filtrage & Enquête Emploi \\
3 & Nettoyage des données & Salaires \\
4 & Restructuration et fusion & Panel Eurostat \\
5 & Statistiques descriptives & Distribution des revenus \\
\midrule
6 & Visualisation & Courbe de Beveridge \\
7 & Régression en pratique I & Équation de Mincer \\
8 & Régression en pratique II & Prix de l'immobilier \\
\midrule
9 & Séries temporelles et API & Rendements boursiers \\
10 & Prédiction et synthèse & Sujet d'annale \\
\bottomrule
\end{tabular}
\end{frame}

\begin{frame}{Économie du travail, macro, micro, finance}
Chaque séance de TD porte sur des données économiques réelles, issues de sources publiques que vous retrouverez en M1.

\vspace{1.2em}
\begin{itemize}
  \item \textbf{INSEE} --- Enquête Emploi, FiLoSoFi, DADS
  \item \textbf{DARES} --- offres et demandes d'emploi
  \item \textbf{DVF} --- transactions immobilières
  \item \textbf{Eurostat}, \textbf{FRED} --- comptes nationaux, séries macroéconomiques
  \item Données de marché --- cours et indices boursiers
\end{itemize}

\vspace{1.2em}
Aucun jeu de données d'exercice. Aucun fichier propre au départ.
\end{frame}

%=========================================================================
\section{Comment cela se passe}
%=========================================================================

\begin{frame}{Cours et travaux dirigés}
\begin{columns}[T]
\begin{column}{0.46\textwidth}
\textbf{En cours}

\vspace{0.5em}
Les concepts, les commandes, et surtout les décisions à prendre.
\end{column}
\begin{column}{0.46\textwidth}
\textbf{En TD}

\vspace{0.5em}
25 min de retour sur l'exercice \\
30 min de résolution en direct \\
45 min de travail autonome
\end{column}
\end{columns}

\vspace{1.8em}
Trois fichiers par séance, sur le site du cours :

\vspace{0.6em}
\begin{itemize}
  \item le \textbf{notebook à trous} --- à ouvrir \emph{avant} le TD ;
  \item le \textbf{script complet} --- mis en ligne après ;
  \item l'\textbf{exercice} --- à traiter pour la fois suivante.
\end{itemize}
\end{frame}

\begin{frame}{Évaluation}
\begin{center}
\begin{tabular}{lccl}
\toprule
\textbf{Épreuve} & \textbf{Poids} & \textbf{Durée} & \textbf{Périmètre} \\
\midrule
Contrôle continu & 40\,\% & 1 h & Séances 1 à 5 \\
Contrôle terminal & 60\,\% & 2 h & Ensemble du cours \\
\bottomrule
\end{tabular}
\end{center}

\vspace{1.5em}
Les deux épreuves se déroulent \alert{sur papier, sans machine}.

\vspace{0.8em}
Aucune ne demande d'écrire du code de mémoire.

\vspace{1.2em}
Les exercices hebdomadaires ne sont pas notés --- mais les questions de diagnostic des épreuves sont construites à partir des erreurs qui y reviennent.
\end{frame}

\begin{frame}{Quatre types de questions}
\begin{enumerate}
  \item \textbf{Lecture de sortie.} Un tableau ou une régression vous est donné : que montre-t-il, et que ne montre-t-il pas ?
  \item \textbf{Diagnostic.} Un extrait de données ou de code contient un problème : l'identifier, proposer une correction.
  \item \textbf{Interprétation économique.} Que peut-on affirmer à partir de ce coefficient, et sous quelles conditions ?
  \item \textbf{Code court.} Trois à cinq lignes à écrire ou compléter --- jamais un script entier.
\end{enumerate}
\vfill
Ce qui est évalué n'est pas la syntaxe. C'est le \alert{jugement sur les données}.
\end{frame}

\begin{frame}{Les assistants d'IA}
Autorisés et attendus pour les exercices. C'est l'outil de travail réel de votre future profession.

Interdits lors des deux épreuves, qui se tiennent sans machine.

\vspace{1.2em}
\begin{block}{Ce qu'ils font bien, ce qu'ils font mal}
Un assistant écrit très bien la syntaxe.

\vspace{0.4em}
Il est peu fiable sur les décisions qui font l'objet de ce cours : quelle variable retenir, comment traiter les manquants, quelle spécification est défendable, ce qu'un coefficient autorise à conclure.
\end{block}

\vspace{0.8em}
Si vous ne pouvez pas expliquer une ligne qu'il a produite, vous ne saurez pas répondre à la question de diagnostic qui portera dessus.
\end{frame}

%=========================================================================
\section{Pour commencer}
%=========================================================================

\begin{frame}{Avant la séance 1}
\begin{enumerate}
  \item Créez un compte Google si vous n'en avez pas.
  \item Ouvrez \texttt{colab.research.google.com} et vérifiez que la page se charge.
  \item Aucune installation n'est nécessaire. Ne téléchargez rien.
\end{enumerate}

\vspace{2em}
Tout le matériel du cours --- calendrier, notebooks, diapositives, syllabus --- se trouve sur :

\vspace{1em}
\begin{center}
\large \texttt{stefaniamarcassa.github.io/analyse\_des\_donnees}
\end{center}
\end{frame}

\begin{frame}[standout]
Des questions ?
\end{frame}

\end{document}
